\section{Real powers along a rectangular intertwiner}
\label{sec:real_power_intertwine}

For a positive semidefinite matrix $A$ and $s\in\R$, define $A^s$
by the spectral calculus. A positive eigenvalue $x$ contributes $x^s$.
At a zero eigenvalue the coefficient is $1$ when $s=0$ and $0$ when
$s\neq0$. Thus a negative power vanishes on the kernel; it does not
assert the existence of an inverse there.
The following auxiliary identity applies to the functional-calculus
calculations in Section~6
of~\cite{OpenAI2026AreaLaw}.

\begin{theorem}[Intertwining real powers]
  \label{thm:real_power_intertwine}
  \lean{Matrix.PosSemidef.rpow_intertwine}
  \leanok
  Let $A$ and $B$ be positive semidefinite matrices on finite-dimensional
  complex Hilbert spaces, and let $J$ be a linear map from the space of
  $B$ to the space of $A$. If $AJ=JB$, then, for every $s\in\R$,
  \begin{align}
    A^sJ&=JB^s.
  \end{align}
\end{theorem}

\begin{proof}
  \leanok
  Choose orthonormal eigenbases of $A$ and $B$, with eigenvalues
  $(a_i)$ and $(b_j)$, and write $K$ for the matrix of $J$ in these
  bases. The intertwining equation gives
  \begin{align}
    (a_i-b_j)K_{ij}&=0.
  \end{align}
  Hence $K_{ij}=0$ whenever $a_i\neq b_j$. Applying the same scalar
  power function, with the stated convention at zero, to both sets of
  eigenvalues gives
  \begin{align}
    a_i^sK_{ij}&=K_{ij}b_j^s.
  \end{align}
  Returning to the original bases proves the assertion. The argument
  requires no isometry condition on $J$ and includes empty index sets.
\end{proof}
